\documentclass[10pt,a4paper]{amsart}
\usepackage[margin=19mm]{geometry}
\usepackage{amsmath,amssymb,mathtools}
\usepackage[scr=rsfs,cal=euler]{mathalfa}
\usepackage{xfrac}
\usepackage[unicode,hidelinks,bookmarks=false]{hyperref}
\newcommand{\ie}{i.e.\ }
\newcommand{\cf}{cf.\ }
\newcommand{\resp}{respectively\ }
\newcommand{\Subsection}{\subsection}
\providecommand{\nobreakdash}{\nobreak\hbox{-}}
\setlength{\emergencystretch}{2em}
\multlinegap0pt
\usepackage{bm}
\usepackage{bbm}
\usepackage{dsfont}
\usepackage{xspace}
\usepackage{cancel}
\usepackage{mathtools}
\usepackage{amsthm}
\newtheorem{theorem}{Theorem}
\newtheorem{lemma}[theorem]{Lemma}
\title[Finite Spectra of Syllogistic Logic with Cardinality Compari]{Finite Spectra of Syllogistic Logic with Cardinality Comparisons}
\author{Ruiting Jiang}
\date{Source: arXiv:2609.24902v1 (CC BY 4.0)}
\begin{document}
\maketitle

\noindent\emph{Source and licence.} Finite Spectra of Syllogistic Logic with Cardinality Comparisons. Version arXiv:2609.24902v1 (CC BY 4.0), distributed under the Creative Commons Attribution 4.0 International licence, as recorded in the arXiv licence metadata for this exact version and shown on the abstract page https://arxiv.org/abs/2609.24902v1. Full rights notice: licenses/arxiv-2609.24902-CC-BY-4.0.txt.\par\smallskip

\noindent\textit{Editorial scope.} Counted unit: the Lemma whose source label is \texttt{eo\_indiscernible} in the article (source0.tex lines 581--583, source label \texttt{eo\_indiscernible}), delivered together with the article's own complete original proof of it (source0.tex lines 584--632, one \texttt{proof} environment of 49 lines). No mathematics is written, repaired or re-derived: statement and proof are the article's own text line for line. Required context retained uncounted: smallest (lines 453--456); o\_indiscernible (lines 495--497); e\_indiscernible (lines 532--534). Every definition, display and label of those lines is retained unchanged. Omitted from this package and not redistributed: 1--452, 457--494, 498--531, 535--580, 633--981. Nothing inside a retained excerpt is omitted or altered: every retained body is delivered exactly as the article's lines are written, so no omission receipt of any kind accompanies this package. Citations: the retained excerpts carry no citation command at all, so no bibliography entry of the article is transcribed and no reference is redistributed. Reference adaptation: none was needed and none was made; every reference inside the delivered unit resolves inside it, so no unresolved reference or citation is produced. Printed slips kept literal: no printed word, symbol, hypothesis or number of the counted statement or of its proof was altered. Layout and class: the article's own class is \texttt{amsart} with packages including amsmath, amsthm, amscd, amsfonts, amssymb, graphicx, xcolor, hyperref; the wrapper is the lane's pinned \texttt{amsart} editorial wrapper at 10pt on A4, which restates the theorem-like environments the retained text needs and adds the article's own macro definitions that the retained text needs (none), with extra wrapper packages (bm, bbm, dsfont, xspace, cancel, mathtools). The delivered numbering is the unit's own. Assets: no retained span contains a figure environment, a graphics-inclusion command, a PDF-page inclusion, a table or a TikZ picture; no image file is copied into this package, the package declares no asset dependency and no image file is redistributed with it. Licence: version arXiv:2609.24902v1 of this article is distributed under the Creative Commons Attribution 4.0 International licence, as recorded in the arXiv licence metadata for this exact version and shown on the abstract page https://arxiv.org/abs/2609.24902v1. A full rights notice is delivered at licenses/arxiv-2609.24902-CC-BY-4.0.txt. Characters: every retained excerpt is pure ASCII, so the delivered engine, XeLaTeX with its default Latin Modern text font, reports no missing character.
\begin{lemma}\label{smallest}
Let $\Gamma$ be any $S^\dagger(card)$-theory. Suppose $n\in Spec(\Gamma)$. Then we have $Spec(\Gamma_n)
\subseteq Spec(\Gamma).$
\end{lemma}

\begin{lemma}\label{o_indiscernible}
	For any odd number $n$,$\{i \mid i \geq n\}\subseteq Spec(\Gamma_n)$.  
\end{lemma}

\begin{lemma}\label{e_indiscernible}
	For any even number $n$, we have $\{2i \mid 2i \geq n \} \subseteq Spec(\Gamma_n)$.
\end{lemma}

\begin{lemma}\label{eo_indiscernible}
	Let $U$ be a spectrum. If an even number $m \in U$, an odd number $n \in U$, and $n > m$, then we have $i \in U$ for all $m \leq i \leq n$. 
\end{lemma}

\begin{proof}
	For an even $i\geq m$, by Lemma~\ref{smallest} we have that $Spec(\Gamma_m)\subseteq U$ and by Lemma~\ref{e_indiscernible} we have $i\in Spec(\Gamma_m).$ Hence $i\in U$.
	
	Thus let $i$ be odd. By Lemma~\ref{smallest} and~\ref{o_indiscernible}, it suffices to prove $m +1 \in U$. 
	
	Let $U = Spec(\Gamma)$. Let $I_0$ be such that $\langle m, I_0\rangle \models \Gamma$ and $I_1$ be such that $\langle n, I_1\rangle \models \Gamma$. Let $P$ be the set of nouns occurring in $\Gamma$, together with their complements. Let $P = Q_0 \cup Q_1$ where $Q_0 = \{x \in P \mid |I_1(x)| < \frac{1}{2}n\}$ and $Q_1 = P \setminus Q_0 = \{x\in P:|I_1(x)|> \frac{1}{2}n\}$. 
	
	Since $n$ is odd, no noun has interpretation of cardinality exactly $\frac{1}{2}n$. Moreover, $|I_1(x)|+|I_1(\bar x)|=n.$ Hence exactly one of $x$ and $\bar x$ belongs to $Q_0$, and the other belongs to $Q_1$. Therefore, renaming variables if necessary, let $Q_0 = \{x_0, \dots, x_l, \dots\}$ and $Q_1 = \{\bar{x_0}, \dots, \bar{x_l}, \dots\}$. 
	
	We are going to add one additional element $m$ to the set $m = \{0, \dots, m-1\}$ and modify the interpretation $I_0$ so that for each raw variable $x$, we either give $m$ to the interpretation of $x$ or to that of $\bar{x}$. Things may go wrong if adding this number changes the cardinality order, or make disjoint sets intersect, so some caution is necessary. 
	
	Define $I: P \to \mathcal{P}(m+1)$ as follows: 
	\[
	I(z) = 
	\begin{cases}
    	 I_0(z),& \mbox{if } z \in Q_0 \\
    	I_0(z) \cup \{m\},              & \mbox{if } z \in Q_1 \\
	\end{cases}
	\]
	It is routine to check that this is an interpretation function, and here we need the fact that $Q_0$ and $Q_1$ are of these forms: $Q_0 = \{x_0, \dots, x_l, \dots\}$ and $Q_1 = \{\bar{x_0}, \dots, \bar{x_l}, \dots\}$. 
	
	We show $\langle m+1, I\rangle \models \Gamma$. Let $\phi \in \Gamma$.	
	\begin{enumerate} 
 		 \item $\phi = \forall(x, y)$ for some $x, y \in P$. 
		 If $x, y \in Q_0$ Then nothing is changed. 	
		 	 
		 If $x, y \in Q_1$ this case follows since the new interpretation function adds the same elements to the interpretation of both variables.		 
		 
		If $x \in Q_0$ and $y \in Q_1$. This is clear since the new function adds the only new elements to the interpretation of all variables in $Q_1$. 		 
		 
		 
		If $x \in Q_1$ and $y \in Q_0$. This is impossible since $\langle n, I_1\rangle \models \forall(x, y)$ we have $I_1(x) \subseteq I_1(y)$ and hence $|I_1(x)| \leq |I_1(y)|$, contradicting $|I_1(x)| > \frac{n}{2} > |I_1(y)|$.
		 
		 Thus, in all possible cases, we have $\langle m+1, I\rangle \models \phi$. 
		 \item $\phi = \exists(x, y)$ for some $x, y \in P$. This is clear because the new interpretation function only adds elements to the interpretation of some variables, so non-empty intersections remain non-empty. 
 		 \item $\phi = \exists^{\geq}(x, y)$ for some $x, y \in P$. 
		 
		  If $x, y \in Q_0$. Then nothing is changed. 
		  
		 If $x, y \in Q_1$. One new element is added to both interpretations, so the cardinality comparison is preserved. 
		 
		 If $x \in Q_0$ and $y \in Q_1$ This is impossible because $|I_1(x)|<n/2<|I_1(y)|$, but $\langle n,I_1\rangle\models\exists^{\geq}(x,y)$ would require $|I_1(x)|\geq|I_1(y)|$.

		 If $x \in Q_1$ and $y \in Q_0$ This still holds because the new interpretation function adds the new element only to the interpretation of $x$. 
  		 \item $\phi = \exists^{>}(x, y)$ for some $x, y \in P$. The same case analysis applies. The case $x\in Q_0$ and $y\in Q_1$ is impossible in the $n$-element model; adding one new element to both sides preserves a strict inequality; and adding it only to the interpretation of $x$ also preserves strict inequality. Hence $|I(x)|>|I(y)|$. 
 	\end{enumerate}

	Therefore $\langle m+1, I\rangle \models \Gamma$, so $m + 1 \in Spec(\Gamma) = U$, which is what remained to be shown.
\end{proof}

\end{document}
